\section{Theory: decisions as tests of theories}\label{sec:theory}

\subsection{The scientific approach to strategy}
Our target behaviour comes from the scientific approach to entrepreneurial and strategic decisions
\citep{camuffo2020scientific,camuffo2024theory}. A decision maker facing an uncertain opportunity should (i) frame the
decision and its real alternatives, including doing nothing; (ii) state \emph{theories}: causal accounts of who would
change their behaviour and why; (iii) derive predictions that tell the theories apart; (iv) gather evidence for and
against each, looking hardest for evidence against the favoured one; (v) when information is worth its cost, run a test
with a decision rule fixed in advance; and (vi) revise when new evidence arrives. Field experiments with start-ups found
that founders trained this way made more accurate decisions, pivoting or terminating sooner when evidence contradicted
them \citep{camuffo2020scientific}.

\subsection{A theory as a restriction on beliefs}
\citet{camuffo2024theory} formalise a theory as a Bayesian network with a restriction. We follow their running example,
Luxottica deciding whether eyeglasses could become a fashion product.

\paragraph{Attributes and states.} An \emph{attribute} is a question with an uncertain answer. Two attributes define the
state space: can the firm design fashion eyewear ($Y$ or $N$), and will demand be high or low ($H$ or $L$). The four
states are $\{YH, YL, NH, NL\}$, and the decision maker cares about the probability of high demand,
\begin{equation}
  P(H) = \theta_Y\,\theta_{H|Y} + (1-\theta_Y)\,\theta_{H|N},
  \label{eq:ph}
\end{equation}
where $\theta_Y = P(Y)$, $\theta_{H|Y} = P(H\mid Y)$ and $\theta_{H|N} = P(H\mid N)$.

\paragraph{The theory.} The theory ``fashion design causes demand'' is one arrow $Y \to H$ in a Bayesian network together
with an inequality on its conditional probabilities,
\begin{equation}
  \Theta = \{\theta : \theta_{H|Y} > \theta_{H|N}\},
  \qquad
  \tilde\Theta = \{\theta : \theta_{H|Y} = \theta_{H|N}\},
\end{equation}
where $\tilde\Theta$ is the null hypothesis that design has no effect.

\paragraph{Value of a theory.} The parameters are unknown, so the decision maker holds a distribution $\mu$ over them. Let
$V_\Theta = \E_{\mu}[P(H)\mid\theta\in\Theta]$ and $V_{\tilde\Theta} = \E_{\mu}[P(H)\mid\theta\in\tilde\Theta]$. With prior
$\omega$ that the theory is right, its value is
\begin{equation}
  V = \omega\,V_\Theta + (1-\omega)\,V_{\tilde\Theta}.
  \label{eq:value}
\end{equation}

\paragraph{A worked example.} Take $\theta_{H|N}=0.5$ and $\theta_Y=0.8$ with probability one, and $\theta_{H|Y}$ uniform.
Under the theory $\theta_{H|Y}$ is uniform on $(0.5,1)$, so $\E[\theta_{H|Y}\mid\Theta]=0.75$ and
\begin{equation}
  V_\Theta = 0.8\cdot 0.75 + 0.2\cdot 0.5 = 0.70 .
\end{equation}
Without the restriction $\theta_{H|Y}$ averages $0.5$ and $P(H)=0.5$. The theory raises belief in high demand by
\emph{ruling parameter values out}: it is informative precisely because it can be wrong.

\paragraph{Experiments.} An experiment produces evidence that updates $\mu$, and with it $V_\Theta$ and
$V_{\tilde\Theta}$. In their account the prior on the theory itself changes only when the evidence moves the beliefs about
parameters far enough, $\lVert \mu' - \mu \rVert > \mu^*$: a small surprise adjusts parameters, a large one shakes belief in
the theory. A rival theory has its own value $Q$, built the same way. With costly experiments the decision maker should
commit once $V$ is far enough above or below $Q$, and experiment in between, the logic of optimal stopping for
costly information \citep{moscarini2001optimal}. Experiments are assumed Bayes plausible: before running one there is no
reason to expect it to move the value in a particular direction.

\begin{remark}[Bayesian updating, illustrated]\label{rem:bayes}
For intuition, suppose $n$ customers are shown the fashion design and $k$ buy. Under the theory, with
$\theta_{H|Y}\sim U(\theta_{H|N},1)$, and under the null, with $\theta_{H|Y}=\theta_{H|N}$, the marginal likelihoods are
\begin{equation}
  m_\Theta(k) = \frac{1}{1-\theta_{H|N}}\int_{\theta_{H|N}}^{1}\binom{n}{k}p^k(1-p)^{n-k}\,dp,
  \qquad
  m_{\tilde\Theta}(k) = \binom{n}{k}\theta_{H|N}^k(1-\theta_{H|N})^{n-k},
\end{equation}
and Bayes' rule gives the posterior on the theory
\begin{equation}
  \omega' = \frac{\omega\, m_\Theta(k)}{\omega\, m_\Theta(k) + (1-\omega)\, m_{\tilde\Theta}(k)}.
\end{equation}
Full updating moves $\omega$ after every observation; the threshold rule of \citet{camuffo2024theory} moves it only after
large surprises, so this calculation overstates how often a theory itself is revised.
\end{remark}

\subsection{From the framework to an agent}
The agent in this paper does \emph{not} compute these quantities, and neither does its environment. A computed update would
need an invented network and invented likelihoods, and assumed ratios could manufacture exactly the revision we want to
measure. Instead the agent carries out the framework's steps in words, with tools and quoted evidence, and the environment
enforces their form:
\begin{itemize}
  \item \emph{theories} are explicit causal claims (who benefits, why behaviour changes, how value is captured), and two of
        them must be genuinely distinct (\cref{sec:env});
  \item \emph{evidence} is a quoted passage from a source, assessed as supporting, contradicting or having no bearing on a
        theory, with what it does not show;
  \item \emph{tests} carry a prediction per theory and a decision rule, within a budget and a time limit;
  \item \emph{beliefs} are stated on a five-step verbal scale (very unlikely, unlikely, uncertain, likely, very likely)
        when the agent takes a position and again after new evidence. The scale is a within-run signal, not a calibrated
        probability; a direction check asks only whether undermining evidence did not raise the affected theory and
        supporting evidence did not lower it.
\end{itemize}
The framework thus supplies the \emph{grammar} of a good investigation; the environment makes that grammar executable;
and fine-tuning teaches a small model to speak it.
